\documentclass[12pt]{exam}
\usepackage[utf8]{inputenc}

\usepackage{amsmath,amstext,amsthm,amssymb,amsxtra}
\usepackage[top=1.5in, bottom=1.5in, left=1.25in, right=1.25in]	{geometry}
%\usepackage[normalem]{ulem}
\usepackage{txfonts} % pxfonts txfonts 
\usepackage[T1]{fontenc}
\usepackage{lmodern}
\renewcommand*\familydefault{\sfdefault}
 \usepackage{euler}   % better than the option below
\usepackage{pdfsync}
\usepackage{multicol}
\newcommand{\ci}[1]{_{ {}_{\scriptstyle #1}}}
\usepackage{graphicx}
\graphicspath{ {images/} }


\newcommand{\norm}[1]{\ensuremath{\left\|#1\right\|}}
\newcommand{\abs}[1]{\ensuremath{\left\vert#1\right\vert}}
\newcommand{\ip}[2]{\ensuremath{\left\langle#1,#2\right\rangle}}
\newcommand{\p}{\ensuremath{\partial}}
\newcommand{\pr}{\mathcal{P}}

\newcommand{\pbar}{\ensuremath{\bar{\partial}}}
\newcommand{\db}{\overline\partial}
\newcommand{\D}{\mathbb{D}}
\newcommand{\B}{\mathbb{B}}
\newcommand{\Sp}{\mathbb{S}}
\newcommand{\T}{\mathbb{T}}
\newcommand{\R}{\mathbb{R}}
\newcommand{\Z}{\mathbb{Z}}
\newcommand{\C}{\mathbb{C}}
\newcommand{\N}{\mathbb{N}}
\newcommand{\Q}{\mathbb{Q}}
\newcommand{\mQ}{\mathcal{Q}}
\newcommand{\mS}{\mathcal{S}}
\newcommand{\scrH}{\mathcal{H}}
\newcommand{\scrL}{\mathcal{L}}
\newcommand{\td}{\widetilde\Delta}
\newcommand{\pw}{\text{PW}}
\newcommand{\esup}{\text{ess.sup}}
\newcommand{\Tn}{\mathcal{T}_n}
\newcommand{\Bn}{\mathbb{B}_n}
\newcommand{\rt}{\mathcal{O}}
\newcommand{\avg}[1]{\langle #1 \rangle}
\newcommand{\one}{\mathbbm{1}}
\newcommand{\eps}{\varepsilon}
\newcommand{\grad}{\nabla}

\newcommand{\La}{\langle }
\newcommand{\Ra}{\rangle }
\newcommand{\rk}{\operatorname{rk}}
\newcommand{\card}{\operatorname{card}}
\newcommand{\ran}{\operatorname{Ran}}
\newcommand{\osc}{\operatorname{OSC}}
\newcommand{\im}{\operatorname{Im}}
\newcommand{\re}{\operatorname{Re}}
\newcommand{\tr}{\operatorname{tr}}
\newcommand{\vf}{\varphi}
\newcommand{\f}[2]{\ensuremath{\frac{#1}{#2}}}

\newcommand{\kzp}{k_z^{(p,\alpha)}}
\newcommand{\klp}{k_{\lambda_i}^{(p,\alpha)}}
\newcommand{\TTp}{\mathcal{T}_p}
\newcommand{\m}[1]{\mathcal{#1}}
\newcommand{\md}{\mathcal{D}}
\newcommand{\qan}{\abs{Q}^{\alpha/n}}
\newcommand{\sbump}[2]{[[ #1,#2 ]]}
\newcommand{\mbump}[2]{\lceil #1,#2 \rceil}
\newcommand{\cbump}[2]{\lfloor #1,#2 \rfloor}

\newcommand{\class}{Math 629 Spring 26}
\newcommand{\term}{Fall 24}
\newcommand{\examnum}{Homework \hwnum}
\newcommand{\examdate}{\duedate}
\newcommand{\timelimit}{75 Minutes}
\newcommand{\vc}[3]{\langle #1,#2,#3\rangle}
\newcommand*{\vv}[1]{\vec{\mkern0mu#1}}
\newcommand{\bv}[1]{\boldsymbol{#1}}
\newcommand{\hide}[1]{}
\newcommand{\uvec}[1]{\boldsymbol{\hat{\textbf{#1}}}}
\newcommand{\vex}[1]{\boldsymbol{{\textbf{#1}}}}

\pagestyle{head}
\firstpageheader{}{}{}
\runningheader{\class}{\examnum\ - Page \thepage\ of \numpages}{\examdate}
\runningheadrule

\makeatletter
\renewcommand*\env@matrix[1][*\c@MaxMatrixCols c]{%
  \hskip -\arraycolsep
  \let\@ifnextchar\new@ifnextchar
  \array{#1}}
\makeatother

%\printanswers
\begin{document}

\noindent
\begin{tabular*}{\textwidth}{l @{\extracolsep{\fill}} r @{\extracolsep{6pt}} l}
\textbf{\class} & \textbf{Name: Mengxiang Jiang}\\
\end{tabular*}\\
\rule[2ex]{\textwidth}{2pt}
%

 

\newcommand{\duedate}{01/19/2026 at 11:59pm}
\newcommand\hwnum{1}

Turn this assginment in on gradescope. All work must be neat and legible. Turn in 
a ``final draft''. There should be nothing scratched out and your work should flow
generally from left to right and top to bottom. 

\begin{questions}
\question
Compute 15 times 17 using the doubling algorithm. 

 \begin{align*}
  `1 &\quad 17\\
  `2 &\quad 34\\
  `4 &\quad 68\\
  `8 &\quad 136\\
  \hline
  Total:\quad 15 &\quad 255
 \end{align*}

\newpage
\question 
Compute $180\div 12$ using the doubling algorithm.\\
This is equivalent to "multiply 12 by what to get 180?"
  \begin{align*}
    1 &\quad 12'\\
    2 &\quad 24'\\
    4 &\quad 48'\\
    8 &\quad 96'\\
    \hline
    Total:\quad 15 &\quad 180
  \end{align*}
  Thus, $180\div 12 = 15$.

\newpage 
\question 
Solve $\frac{7}{5}x + 3 = 24$ using false position. 
\begin{align*}
  \text{Guess } x=5: &\quad \frac{7}{5}(5) + 3 = 10 \quad (\text{Not high enough})\\
  \text{Guess } x=10: &\quad \frac{7}{5}(10) + 3 = 17 \quad (\text{Still not high enough})\\
  \text{Guess } x=15: &\quad \frac{7}{5}(15) + 3 = 24 \quad (\text{Correct})\\
\end{align*}

\newpage 
\question 
Solve $\frac{6}{5}x + 3 = 20$ using false position. 
\begin{align*}
  \text{Guess } x=5: &\quad \frac{6}{5}(5) + 3 = 9 \quad (\text{Not high enough})\\
  \text{Guess } x=10: &\quad \frac{6}{5}(10) + 3 = 15 \quad (\text{Still not high enough})\\
  \text{Guess } x=15: &\quad \frac{6}{5}(15) + 3 = 21 \quad (\text{Too high})\\
  \text{Guess } x=14: &\quad \frac{6}{5}(14) + 3 = 19 + \frac{4}{5} \quad \left(\text{Not high enough, need } \frac{1}{5}\right)\\
  \text{Guess } x=14 + \frac{1}{6}: &\quad \frac{6}{5}\left(14 + \frac{1}{6}\right) + 3 = 20 \quad (\text{Correct})\\
\end{align*}

\end{questions}
\end{document}
